\documentclass{article}
\usepackage{amsmath, amssymb, amsthm}
\usepackage{hyperref}
\usepackage{geometry}
\geometry{margin=1in}

\title{Erd\H{o}s Problem \#289}
\date{Last updated: 2025-08-31}
\author{Source: erdosproblems.com}

\begin{document}
\maketitle

\noindent\textbf{Status:} OPEN \quad \textbf{No prize}

\noindent\textbf{Tags:} number theory, unit fractions

\medskip
\noindent\textbf{Problem Statement:}

Is it true that, for all sufficiently large $k$, there exist finite intervals $I_1,\ldots,I_k\subset \mathbb{N}$, distinct, not overlapping or adjacent, with $\lvert I_i\rvert \geq 2$ for $1\leq i\leq k$ such that\[1=\sum_{i=1}^k \sum_{n\in I_i}\frac{1}{n}?\]

\bigskip
\noindent\textbf{Remarks:}

Erd\H{o}s and Graham posed this in \cite{ErGr80} without the stipulation the intervals be distinct, non-overlapping, or adjacent, but Kovac in the comments has provided a simple argument showing that it is easily possible without this restriction, and likely \cite{ErGr80} just forgot to mention this natural restriction.

As an example representing $2$ rather than $1$, Hickerson and Montgomery, in the solution to AMS Monthly problem E2689 proposed by Hahn, found\[2=\sum_{i=1}^5 \sum_{n\in I_i}\frac{1}{n}\]where $I_1=[2,7]$, $I_2=[9,10]$, $I_3=[17,18]$, $I_4=[34,35]$, and $I_5=[84,85]$.

\bigskip
\noindent\textbf{References:}

[ErGr80] Erd\H{o}s, P. and Graham, R., Old and new problems and results in combinatorial number theory. Monographies de L'Enseignement Mathematique (1980).

\bigskip
\noindent\small{Source: \url{https://www.erdosproblems.com/289}}
\end{document}
